% Chapter 3 - C.

\guidechapter{PROGRAMMING IN C}{%
  \item Building against the SDK
  \item The shape of every call
  \item Buffers belong to you
  \item Result codes
  \item Capability detection
  \item A worked example
}

\section{Building Against The SDK}

Build the library once:

\begin{listingbox}
make lib
\end{listingbox}

then compile against \cmd{include/} and link \cmd{ultimate.lib}:

\begin{listingbox}
cl65 -t c64 -Osir -I include -o prog.prg prog.c \textbackslash\\
\hspace*{2em}bindings/cc65/build/ultimate.lib
\end{listingbox}

One header covers the whole library:

\begin{listingbox}
\#include <ultimate.h>
\end{listingbox}

If you need the transport itself --- to send a command the SDK does not wrap ---
add \cmd{<uci.h>} as well. Most programs never do.

\section{The Shape Of Every Call}

Almost every function in the SDK looks like this:

\begin{listingbox}
uint8\_t ultimate\_something(\ldots);
\end{listingbox}

It returns a result code, and \cmd{ULTIMATE\_OK} is zero. Anything it produces
comes back through a pointer you passed in. There are no global error variables
to check, nothing is returned in a static buffer, and no call allocates memory.

\begin{listingbox}
uint16\_t got;\\
uint8\_t  err;\\
\\
err = ultimate\_read(buf, sizeof buf, \&got);\\
if (err != ULTIMATE\_OK)\\
\hspace*{1em}return err;
\end{listingbox}

\section{Buffers Belong To You}

The SDK has no heap and no hidden buffers. Every byte it produces lands in
memory you own, and it never keeps a pointer after the call returns. That is
what makes it usable from an interrupt-driven demo, and it is why every call
that produces data takes both a buffer and its size.

\begin{warning}
The size you pass is the size of the buffer, not the number of bytes you want.
If a reply is larger, you get \cmd{ULTIMATE\_ERR\_TRUNCATED}, the count tells
you how much was kept, and nothing is written past the end.
\end{warning}

\section{Result Codes}

There are eleven, they never change, and Appendix~B lists them. Three are worth
knowing straight away:

\begin{reftable}{lp{0.60\linewidth}}
\cmd{ULTIMATE\_OK} & it worked \\
\cmd{ULTIMATE\_END} & there is no more --- the end of a directory, or a socket
  the far end has closed. Not a failure \\
\cmd{ULTIMATE\_ERR\_TRUNCATED} & it worked, and your buffer was too small \\
\end{reftable}

Turn any of them into text with:

\begin{listingbox}
printf("\%s\textbackslash n", ultimate\_strerror(err));
\end{listingbox}

\section{Capability Detection}

Different machines have different targets, and firmware versions differ too.
Ask once, at start-up, rather than assuming:

\begin{listingbox}
ultimate\_capabilities caps;\\
\\
ultimate\_detect(\&caps);\\
if (ultimate\_has\_http(\&caps)) \{ /* fetch things */ \}\\
if (ultimate\_has\_network(\&caps)) \{ /* sockets */ \}
\end{listingbox}

The hardware that is not reached through the command interface answers
separately, because it is not a target: \cmd{ultimate\_turbo\_available()} and
\cmd{ultimate\_reu\_size()} tell you about the processor speed registers and
the RAM expansion.

\section{A Worked Example}

This prints the directory, one entry at a time. It is the pattern for every
walk in the SDK: start it, take entries until \cmd{ULTIMATE\_END}, and check
the result of the call that ended it.

\begin{listingbox}
\#include <stdio.h>\\
\#include <ultimate.h>\\
\\
int main(void)\\
\{\\
\hspace*{1em}char    name[64];\\
\hspace*{1em}uint8\_t attrib;\\
\hspace*{1em}uint8\_t err;\\
\\
\hspace*{1em}if (ultimate\_init() != ULTIMATE\_OK) \{\\
\hspace*{2em}printf("no ultimate\textbackslash n");\\
\hspace*{2em}return 1;\\
\hspace*{1em}\}\\
\\
\hspace*{1em}err = ultimate\_opendir();\\
\hspace*{1em}while (err == ULTIMATE\_OK) \{\\
\hspace*{2em}err = ultimate\_readdir(name, sizeof name, \&attrib);\\
\hspace*{2em}if (err != ULTIMATE\_OK)\\
\hspace*{3em}break;\\
\hspace*{2em}printf("\%s\%s\textbackslash n", name,\\
\hspace*{3em}(attrib \& DOS\_ATTR\_DIR) ? "/" : "");\\
\hspace*{1em}\}\\
\\
\hspace*{1em}if (err != ULTIMATE\_END)\\
\hspace*{2em}printf("stopped: \%s\textbackslash n", ultimate\_strerror(err));\\
\hspace*{1em}return 0;\\
\}
\end{listingbox}

\begin{note}
A directory walk is one live exchange with the Ultimate. Do not issue any other
command in the middle of one --- the block boundaries between entries are the
only thing separating one name from the next. If you need to stop early, call
\cmd{uci\_abort()}.
\end{note}
